% Read STYLE.md before editing prose in this file.
% Follow the Content Creation Workflow in WORKFLOW.md §Content Creation Workflow.

\section{Methodology}
\label{sec:methodology}

\subsection{Data Collection}
\label{subsec:data-collection}

Acoustic recordings were obtained at three Atlantic sites: Site~A
(western Azores archipelago), Site~B (Bay of Biscay, northern Spain),
and Site~C (Cape Verde islands).
Field observations were conducted between June~2021 and September~2024,
yielding 42~confirmed departure events for which full acoustic records
were available.
A departure event was defined as a session in which a marked individual
or coherent group left the observation area and did not return within
the 72-hour monitoring window.

\begin{definition}[Departure Event]
	\label{def:departure-event}
	A \emph{departure event} is a continuous acoustic recording window of
	\(\pm 5\) minutes around the last observed surfacing of a focal
	individual or group, provided that individual or group did not re-enter
	the monitoring area within 72~hours of the recording.
\end{definition}

Recordings were made using calibrated hydrophone arrays deployed at
depths of 3--8~m, sampling at 96\kHz{} with 24-bit resolution.
All recordings were stored in lossless FLAC format and are archived
with the \citet{GuideResearch2025dataset} dataset.

\subsection{Annotation Scheme}
\label{subsec:annotation}

Each departure event was segmented into discrete acoustic units by two
independent annotators following the protocol in
\citet{GuideResearch2025dataset}.
Inter-annotator agreement, measured by Cohen's~$\kappa$, was $0.83$
across the full corpus---within the range typically considered
``substantial'' agreement~\citep{Dent2024farewell}.

Each unit was assigned one of five labels:
\texttt{whistle}, \texttt{burst-pulse}, \texttt{click-train},
\texttt{silence}, or \texttt{noise}.
The sequence of labels for a departure event constitutes a
\emph{vocalization trace}.

\subsection{Computational Analysis}
\label{subsec:analysis}

\paragraph{Hidden Markov Model}

We model each vocalization trace as a sequence emitted by a discrete-time
HMM with $N$ hidden states representing acoustic context:

\begin{equation}
	\label{eq:hmm-likelihood}
	P(\mathbf{o} \mid \lambda) =
	\sum_{\mathbf{q}} P(\mathbf{o} \mid \mathbf{q}, \lambda)\,
	P(\mathbf{q} \mid \lambda),
\end{equation}

where $\mathbf{o} = (o_1, \ldots, o_T)$ is the observed label sequence,
$\mathbf{q} = (q_1, \ldots, q_T)$ is the hidden state sequence, and
$\lambda = (A, B, \pi)$ denotes the transition matrix, emission matrix,
and initial distribution, respectively.

Model parameters were estimated using the Baum--Welch algorithm with
random restarts ($k=50$).
The number of hidden states $N$ was selected by Bayesian information
criterion over the range $N \in \{2, 3, \ldots, 10\}$; the
optimal value was $N = 4$.

\paragraph{Spectral Envelope Analysis}

For each annotated whistle unit, we extracted the spectral envelope using
linear predictive coding (LPC) with order $p = 16$.
The resulting envelope coefficients were projected onto a 2D space via
principal component analysis (PCA) for visualization.

\Cref{tab:vocalization-comparison} summarizes the acoustic features of
the three identified archetypes.

%% tables/comparison_table.tex — Vocalization archetype comparison table
%% Included via %% tables/comparison_table.tex — Vocalization archetype comparison table
%% Included via %% tables/comparison_table.tex — Vocalization archetype comparison table
%% Included via \input{tables/comparison_table} in chapters/03_methodology.tex

\begin{table}[htbp]
  \caption{Acoustic features of the three departure vocalization archetypes
    identified in the Atlantic Cetacean Farewell Corpus (ACFC v1.0).
    Values are mean $\pm$ standard deviation across all instances.
    Frequency ranges are for the dominant whistle component.}
  \label{tab:vocalization-comparison}
  \begin{tabularx}{\linewidth}{>{{\raggedright\arraybackslash}}X r r r >{\raggedright\arraybackslash}X}
    \toprule
    \textbf{Archetype}
      & \makecell[r]{\textbf{Peak Freq.}\\\textbf{(\kHz)}}
      & \textbf{Duration (s)}
      & \textbf{Count}
      & \textbf{Characteristic Feature} \\
    \midrule
    \aseq{brief-farewell}
      & $14.2 \pm 1.8$
      & $0.83 \pm 0.21$
      & 18
      & Single upward frequency sweep; minimal burst-pulse content \\[3pt]
    \aseq{extended-farewell}
      & $11.6 \pm 2.3$
      & $3.47 \pm 0.94$
      & 17
      & Multi-phrase structure with recurring motif; broad spectral envelope \\[3pt]
    \aseq{urgent-departure}
      & $17.1 \pm 0.9$
      & $0.41 \pm 0.11$
      & 7
      & High-frequency, short-duration; high burst-pulse density \\
    \bottomrule
  \end{tabularx}
\end{table}
 in chapters/03_methodology.tex

\begin{table}[htbp]
  \caption{Acoustic features of the three departure vocalization archetypes
    identified in the Atlantic Cetacean Farewell Corpus (ACFC v1.0).
    Values are mean $\pm$ standard deviation across all instances.
    Frequency ranges are for the dominant whistle component.}
  \label{tab:vocalization-comparison}
  \begin{tabularx}{\linewidth}{>{{\raggedright\arraybackslash}}X r r r >{\raggedright\arraybackslash}X}
    \toprule
    \textbf{Archetype}
      & \makecell[r]{\textbf{Peak Freq.}\\\textbf{(\kHz)}}
      & \textbf{Duration (s)}
      & \textbf{Count}
      & \textbf{Characteristic Feature} \\
    \midrule
    \aseq{brief-farewell}
      & $14.2 \pm 1.8$
      & $0.83 \pm 0.21$
      & 18
      & Single upward frequency sweep; minimal burst-pulse content \\[3pt]
    \aseq{extended-farewell}
      & $11.6 \pm 2.3$
      & $3.47 \pm 0.94$
      & 17
      & Multi-phrase structure with recurring motif; broad spectral envelope \\[3pt]
    \aseq{urgent-departure}
      & $17.1 \pm 0.9$
      & $0.41 \pm 0.11$
      & 7
      & High-frequency, short-duration; high burst-pulse density \\
    \bottomrule
  \end{tabularx}
\end{table}
 in chapters/03_methodology.tex

\begin{table}[htbp]
  \caption{Acoustic features of the three departure vocalization archetypes
    identified in the Atlantic Cetacean Farewell Corpus (ACFC v1.0).
    Values are mean $\pm$ standard deviation across all instances.
    Frequency ranges are for the dominant whistle component.}
  \label{tab:vocalization-comparison}
  \begin{tabularx}{\linewidth}{>{{\raggedright\arraybackslash}}X r r r >{\raggedright\arraybackslash}X}
    \toprule
    \textbf{Archetype}
      & \makecell[r]{\textbf{Peak Freq.}\\\textbf{(\kHz)}}
      & \textbf{Duration (s)}
      & \textbf{Count}
      & \textbf{Characteristic Feature} \\
    \midrule
    \aseq{brief-farewell}
      & $14.2 \pm 1.8$
      & $0.83 \pm 0.21$
      & 18
      & Single upward frequency sweep; minimal burst-pulse content \\[3pt]
    \aseq{extended-farewell}
      & $11.6 \pm 2.3$
      & $3.47 \pm 0.94$
      & 17
      & Multi-phrase structure with recurring motif; broad spectral envelope \\[3pt]
    \aseq{urgent-departure}
      & $17.1 \pm 0.9$
      & $0.41 \pm 0.11$
      & 7
      & High-frequency, short-duration; high burst-pulse density \\
    \bottomrule
  \end{tabularx}
\end{table}

